%LTeX: language=fr-FR
\documentclass[a4paper,11pt]{article}

% ---------- Préambule minimal ----------
\usepackage[french]{babel}
\usepackage{amsmath,amssymb}
\usepackage{enumitem}
\usepackage{tikz}
\usepackage{tasks}
\usetikzlibrary{arrows.meta}
\usepackage[margin=2.5cm]{geometry}
\usepackage{needspace}
\usepackage{currfile}
\usepackage{fancyhdr}
\usepackage{lastpage}
\usepackage{array}
\usepackage{colortbl}
\usepackage{tcolorbox}
\settasks{label=\arabic*.}

\pagestyle{empty}
\graphicspath{ {../Images/} }
\input{\currfiledir ../def.tex}
\def \Document {Exercices}


% Environnement "exercice" : \begin{exercice}{<numéro>} ... \end{exercice}
\newenvironment{exercice}[1]{%
  \Needspace{6\baselineskip}\par\bigskip\noindent\textbf{Exercice #1}\par\smallskip\noindent\ignorespaces
}{%
  \par
}

% Titre de partie, en lien avec les parties du cours
\newcommand{\partie}[1]{%
  \Needspace{10\baselineskip}\par\bigskip\noindent{\large\textbf{#1}}\par\vspace{-0.6em}\noindent\rule{\linewidth}{0.4pt}\par
}

% Repère pour une suite : \grille{nmax}{ymax}{graduations en y}{pas de la grille en y}
\newcommand{\grille}[4]{%
  \draw[very thin, gray!40, xstep=1, ystep=#4] (0,0) grid (#1,#2);
  \draw[->] (0,0) -- ({#1+0.5},0) node[right] {\small $n$};
  \draw[->] (0,0) -- (0,#2) -- ++(0,0.4cm);
  \foreach \x in {1,...,#1} \node[below] at (\x,0) {\scriptsize $\x$};
  \foreach \y in {#3} \node[left] at (0,\y) {\scriptsize $\y$};
  \node[below left] at (0,0) {\scriptsize $0$};
}

\newcolumntype{C}[1]{>{\centering\arraybackslash}m{#1}}

% Cadre pour les programmes Python
\newtcolorbox{python}{colback=white, colframe=black!50, boxrule=0.5pt, arc=1pt, left=6pt, fontupper=\ttfamily}
\newcommand{\tab}{\hspace*{2em}}

\begin{document}
\pagestyle{fancy}

\fancyhf{} % clear existing header/footer entries
\fancyhead[L]{ {\Niveau}}
\fancyhead[C]{ {\Chapitre}}
\fancyhead[R]{ {\Document}}
\fancyfoot[L]{ A. Fernandez }
\fancyfoot[R]{{\thepage\  / \pageref{LastPage}}}

% =====================================================
\partie{1. Définition, notation $u_n$, formule explicite}

\begin{exercice}{1}
Un centre de vaccination note chaque semaine le nombre de vaccins réalisés. On note $u_n$ le nombre de vaccins de la semaine $n$.

\renewcommand{\arraystretch}{1.4}
\begin{center}
\begin{tabular}{|c|*{6}{C{1.1cm}|}}
  \hline
  $n$ & $0$ & $1$ & $2$ & $3$ & $4$ & $5$ \\
  \hline
  $u_n$ & $120$ & $135$ & $150$ & $142$ & $160$ & $175$ \\
  \hline
\end{tabular}
\end{center}
\begin{enumerate}[label=\arabic*.]
\item Donner le premier terme de la suite.
\item Donner $u_3$ et interpréter ce nombre dans le contexte.
\item Quel est le rang du terme égal à $150$~?
\item Calculer $u_{2+1}$ puis $u_2 + 1$. Ces deux nombres sont-ils égaux~?
\end{enumerate}
\end{exercice}

% =====================================================
\begin{exercice}{2}
Dans chaque cas, calculer les termes demandés.
\begin{tasks}(2)
\task $u_n = 4n - 3$ : $u_0$, $u_1$ et $u_{10}$.
\task $v_n = n^2 - 2n$ : $v_0$, $v_3$ et $v_{10}$.
\task $w_n = \dfrac{60}{n + 1}$ : $w_0$, $w_2$ et $w_5$.
\task $t_n = 100 - 7n$ : $t_0$, $t_4$ et $t_{10}$.
\end{tasks}
\end{exercice}

% =====================================================
\begin{exercice}{3}
Une association distribue des colis alimentaires. On estime que le nombre de colis distribués la semaine $n$ est $c_n = 150 + 12n$.
\begin{enumerate}[label=\arabic*.]
\item Calculer $c_0$ et interpréter ce nombre.
\item Combien de colis seront distribués la semaine $8$~?
\item Compléter le tableau suivant.

\renewcommand{\arraystretch}{1.4}
\begin{tabular}{|c|*{6}{C{1.1cm}|}}
  \hline
  $n$ & $0$ & $1$ & $2$ & $3$ & $4$ & $5$ \\
  \hline
  $c_n$ & & & & & & \\
  \hline
\end{tabular}
\item Quelle semaine l'association distribuera-t-elle $270$ colis~? On pourra résoudre l'équation $150 + 12n = 270$.
\end{enumerate}
\end{exercice}

% =====================================================
\partie{2. Relation de récurrence, comparer les deux modes}

\begin{exercice}{4}
Dans chaque cas, calculer $u_1$, $u_2$, $u_3$ et $u_4$ en détaillant les calculs.
\begin{tasks}(2)
\task $u_0 = 2$ et $u_{n+1} = u_n + 5$
\task $u_0 = 1$ et $u_{n+1} = 3 \times u_n$
\task $u_0 = 10$ et $u_{n+1} = 2 \times u_n - 4$
\task $u_0 = 100$ et $u_{n+1} = 0{,}9 \times u_n$
\end{tasks}
\end{exercice}

% =====================================================
\begin{exercice}{5}
Pour chaque suite, indiquer si elle est définie par une formule explicite ou par une relation de récurrence.
\begin{tasks}(2)
\task $u_n = 7n + 2$
\task $v_0 = 4$ et $v_{n+1} = v_n - 1$
\task $w_n = 3 \times 2^n$
\task $t_0 = 50$ et $t_{n+1} = 1{,}2 \times t_n$
\task $r_n = n^2 + n$
\task $s_0 = 1$ et $s_{n+1} = s_n^2 + 1$
\end{tasks}
Pour quelles suites peut-on calculer directement le terme de rang $20$~? Le calculer pour la suite $u$.
\end{exercice}

% =====================================================
\begin{exercice}{6}
On considère la suite $u$ définie par $u_0 = 1$ et $u_{n+1} = u_n + 3$, et la suite $v$ définie par $v_n = 3n + 1$.
\begin{enumerate}[label=\arabic*.]
\item Calculer $u_1$, $u_2$, $u_3$ et $u_4$.
\item Calculer $v_0$, $v_1$, $v_2$, $v_3$ et $v_4$. Que remarque-t-on~?
\item Calculer $u_{20}$ le plus rapidement possible. Expliquer la méthode.
\end{enumerate}
\end{exercice}

% =====================================================
\begin{exercice}{7}
Dans une pharmacie hospitalière, on note $s_n$ le stock de gel hydroalcoolique, en litres, au début de la semaine $n$. On a $s_0 = 200$ et $s_{n+1} = 0{,}5 \times s_n + 40$.
\begin{enumerate}[label=\arabic*.]
\item Calculer $s_1$, $s_2$, $s_3$ et $s_4$.
\item Expliquer avec une phrase comment le stock évolue d'une semaine à la suivante.
\end{enumerate}
\end{exercice}

% =====================================================
\partie{3. Représentation graphique}

\begin{exercice}{8}
Représenter chaque suite par un nuage de points, pour $n$ allant de $0$ à $5$ (pour $v$, de $0$ à $4$).

\medskip
\noindent
\begin{minipage}[t]{0.48\linewidth}
\centering
$u_n = 10 - 2n$

\smallskip
\begin{tikzpicture}[x=0.75cm, y=0.35cm]
  \grille{6}{11}{2,4,...,10}{1}
\end{tikzpicture}
\end{minipage}
\hfill
\begin{minipage}[t]{0.48\linewidth}
\centering
$v_0 = 1$ et $v_{n+1} = 2 \times v_n$

\smallskip
\begin{tikzpicture}[x=0.75cm, y=0.22cm]
  \grille{6}{17}{2,4,...,16}{1}
\end{tikzpicture}
\end{minipage}

\medskip
Pour quelle suite les points sont-ils alignés~? Pouvait-on le prévoir~?
\end{exercice}

% =====================================================
\begin{exercice}{9}
\begin{minipage}[c]{0.5\linewidth}
On a représenté ci-contre le nombre $\ell_n$ de lits occupés dans un service hospitalier le jour $n$.
\begin{enumerate}[label=\arabic*.]
\item Lire $\ell_0$, $\ell_2$ et $\ell_5$.
\item Quel jour $25$ lits sont-ils occupés~?
\item Quel jour le nombre de lits occupés est-il le plus grand~?
\item Les points sont-ils alignés~?
\end{enumerate}
\end{minipage}
\hfill
\begin{minipage}[c]{0.46\linewidth}
\centering
\begin{tikzpicture}[x=0.75cm, y=0.15cm]
  \grille{7}{30}{5,10,...,30}{1}
  \foreach \n/\u in {0/18, 1/22, 2/25, 3/21, 4/24, 5/28, 6/26} \fill[blue] (\n,\u) circle (3pt);
\end{tikzpicture}
\end{minipage}
\end{exercice}

% =====================================================
\Needspace{16\baselineskip}
\begin{exercice}{10}
Associer chaque suite à son nuage de points, en calculant quelques termes.
\begin{tasks}(3)
\task $u_n = 2n + 1$
\task $v_n = n^2$
\task $w_0 = 16$ et $w_{n+1} = w_n \div 2$
\end{tasks}

\medskip
\noindent
\begin{minipage}[t]{0.32\linewidth}
\centering
Nuage A

\begin{tikzpicture}[x=0.6cm, y=0.15cm]
  \grille{5}{17}{4,8,12,16}{1}
  \foreach \n/\u in {0/16, 1/8, 2/4, 3/2, 4/1} \fill[blue] (\n,\u) circle (2.5pt);
\end{tikzpicture}
\end{minipage}
\hfill
\begin{minipage}[t]{0.32\linewidth}
\centering
Nuage B

\begin{tikzpicture}[x=0.6cm, y=0.15cm]
  \grille{5}{17}{4,8,12,16}{1}
  \foreach \n/\u in {0/1, 1/3, 2/5, 3/7, 4/9} \fill[blue] (\n,\u) circle (2.5pt);
\end{tikzpicture}
\end{minipage}
\hfill
\begin{minipage}[t]{0.32\linewidth}
\centering
Nuage C

\begin{tikzpicture}[x=0.6cm, y=0.15cm]
  \grille{5}{17}{4,8,12,16}{1}
  \foreach \n/\u in {0/0, 1/1, 2/4, 3/9, 4/16} \fill[blue] (\n,\u) circle (2.5pt);
\end{tikzpicture}
\end{minipage}
\end{exercice}

% =====================================================
\partie{4. Modéliser une situation par une suite}

\begin{exercice}{11}
Dans chaque cas, donner le premier terme et la relation de récurrence de la suite, puis calculer les deux termes suivants.
\begin{enumerate}[label=\arabic*.]
\item Une crèche a $30$ places en 2026. Elle en crée $4$ de plus chaque année. On note $p_n$ le nombre de places en $2026 + n$.
\item Une population de $2\,000$ bactéries augmente de $15\,\%$ chaque heure. On note $b_n$ le nombre de bactéries au bout de $n$ heures.
\item Un service dispose de $500$ masques et en utilise $40$ par jour. On note $m_n$ le nombre de masques restants au bout de $n$ jours.
\item Un appareil d'échographie vaut $80\,000$~€ à l'achat. Sa valeur baisse de $10\,\%$ chaque année. On note $e_n$ sa valeur au bout de $n$ années.
\end{enumerate}
\end{exercice}

% =====================================================
\begin{exercice}{12}
Un patient reçoit une perfusion. Au départ, il a $20$~mg de médicament dans le sang. Chaque heure, son organisme élimine $20\,\%$ du médicament, puis la perfusion en apporte $5$~mg. On note $q_n$ la quantité de médicament, en mg, au bout de $n$ heures.
\begin{enumerate}[label=\arabic*.]
\item Justifier que $q_{n+1} = 0{,}8 \times q_n + 5$.
\item Calculer $q_1$, $q_2$, $q_3$ et $q_4$. Arrondir au centième si besoin.
\item À chaque heure, ajoute-t-on toujours la même quantité~? Multiplie-t-on toujours par le même nombre~?
\end{enumerate}
\end{exercice}

% =====================================================
\partie{5. Sens de variation}

\begin{exercice}{13}
Pour chaque suite, conjecturer son sens de variation à partir du tableau de valeurs.

\renewcommand{\arraystretch}{1.3}
\begin{center}
\begin{tabular}{|c|*{6}{C{0.9cm}|}}
  \hline
  $n$ & $0$ & $1$ & $2$ & $3$ & $4$ & $5$ \\
  \hline
  $u_n$ & $3$ & $5$ & $8$ & $12$ & $17$ & $23$ \\
  \hline
  $v_n$ & $50$ & $45$ & $41$ & $38$ & $36$ & $35$ \\
  \hline
  $w_n$ & $4$ & $6$ & $5$ & $7$ & $6$ & $8$ \\
  \hline
\end{tabular}
\end{center}
Quel est le sens de variation de la suite $\ell$ de l'exercice 9~?
\end{exercice}

% =====================================================
\begin{exercice}{14}
Dans chaque cas, calculer $u_{n+1} - u_n$ et en déduire le sens de variation de la suite.
\begin{tasks}(2)
\task $u_0 = 7$ et $u_{n+1} = u_n + 2{,}5$
\task $u_0 = 40$ et $u_{n+1} = u_n - 6$
\task $u_n = 5n - 3$
\task $u_n = 12 - 4n$
\task $u_0 = 3$ et $u_{n+1} = u_n + n$
\task $u_n = 8 + 0{,}5n$
\end{tasks}
\end{exercice}

% =====================================================
\partie{6. Suites, tableur et Python}

\begin{exercice}{15}
\begin{minipage}[c]{0.42\linewidth}
\renewcommand{\arraystretch}{1.25}
\begin{tabular}{|>{\columncolor{gray!20}}c|C{0.8cm}|C{1.6cm}|C{1cm}|}
  \hline
  \rowcolor{gray!20} & A & B & C \\
  \hline
  1 & $n$ & $b_n$ & $p_n$ \\
  \hline
  2 & 0 & 2\,000 & 30 \\
  \hline
  3 & 1 & 2\,300 & 34 \\
  \hline
  4 & 2 & 2\,645 & 38 \\
  \hline
  5 & 3 & 3\,041,75 & 42 \\
  \hline
  6 & 4 & 3\,498,01 & 46 \\
  \hline
  7 & 5 & 4\,022,71 & 50 \\
  \hline
  8 & 6 & 4\,626,12 & 54 \\
  \hline
\end{tabular}
\end{minipage}
\hfill
\begin{minipage}[c]{0.55\linewidth}
On reprend deux suites de l'exercice 11 : $b_0 = 2\,000$ et $b_{n+1} = 1{,}15 \times b_n$ (bactéries), et $p_n = 30 + 4n$ (places en crèche).
\begin{enumerate}[label=\arabic*.]
\item Quelle formule a-t-on saisie en C2, puis recopiée vers le bas~?
\item Quelle formule a-t-on saisie en B3, puis recopiée vers le bas~?
\item Que devient la formule de B3 en B6~?
\item À partir de quelle heure y a-t-il plus de $4\,000$ bactéries~?
\item À partir de quelle année la crèche aura-t-elle au moins $50$ places~?
\end{enumerate}
\end{minipage}
\end{exercice}

% =====================================================
\begin{exercice}{16}
\begin{minipage}[t]{0.4\linewidth}
\vspace{0pt}
\begin{python}
  def u(n):\\
  \tab return 3*n + 2
\end{python}
\begin{python}
  def v(n):\\
  \tab terme = 5\\
  \tab for i in range(n):\\
  \tab\tab terme = terme + 4\\
  \tab return terme
\end{python}
\end{minipage}
\hfill
\begin{minipage}[t]{0.56\linewidth}
\vspace{0pt}
\begin{enumerate}[label=\arabic*.]
\item Que renvoient \texttt{u(0)} et \texttt{u(4)}~? La suite $u$ est-elle définie par une formule explicite ou par récurrence~?
\item Que renvoient \texttt{v(1)} et \texttt{v(3)}~? Donner $v_0$ et la relation de récurrence de la suite $v$.
\item Écrire une fonction \texttt{w(n)} qui renvoie le terme $w_n$ de la suite définie par $w_0 = 1\,000$ et $w_{n+1} = 0{,}8 \times w_n + 50$.
\item Écrire une fonction \texttt{t(n)} qui renvoie $t_n = n^2 + 1$. On pourra écrire $n^2$ sous la forme \texttt{n*n}.
\end{enumerate}
\end{minipage}
\end{exercice}

% =====================================================
\newpage
\partie{7. Synthèse}

\begin{exercice}{17}
Deux centres de soins comptaient des patients réguliers en 2026 :
\begin{itemize}
\item le centre A avait $500$ patients, et en gagne $40$ de plus chaque année ;
\item le centre B avait $400$ patients, et leur nombre augmente de $10\,\%$ chaque année.
\end{itemize}
On note $a_n$ et $b_n$ le nombre de patients des centres A et B en $2026 + n$.
\begin{enumerate}[label=\arabic*.]
\item Donner $a_0$, $b_0$ et les relations de récurrence des suites $a$ et $b$.
\item Justifier que $a_n = 500 + 40n$.
\item Compléter le tableau (arrondir $b_n$ à l'unité).

\renewcommand{\arraystretch}{1.4}
\begin{tabular}{|c|*{9}{C{0.95cm}|}}
  \hline
  $n$ & $0$ & $1$ & $2$ & $3$ & $4$ & $5$ & $6$ & $7$ & $8$ \\
  \hline
  $a_n$ & & & & & & & & & \\
  \hline
  $b_n$ & & & & & & & & & \\
  \hline
\end{tabular}

\item Représenter les deux suites dans le repère ci-dessous. Quel nuage est formé de points alignés~?

\begin{center}
\begin{tikzpicture}[x=1.2cm, y=0.0065cm]
  \grille{8}{900}{100,200,...,900}{50}
\end{tikzpicture}
\end{center}

\item Déterminer le sens de variation de la suite $a$ en calculant $a_{n+1} - a_n$. Conjecturer celui de la suite $b$.
\item En quelle année le centre B aura-t-il plus de patients que le centre A~?
\end{enumerate}
\end{exercice}

\end{document}
